\documentclass[UTF8,12pt,a4paper]{ctexart}

\usepackage{amsmath}
\usepackage{cases}
\usepackage{cite}
\usepackage{graphicx}
\usepackage{enumerate}
\usepackage{algorithm}
\usepackage{caption}   % \caption*需要
\usepackage{subcaption} % 子图布局
\usepackage[noend]{algpseudocode} %algorithmicx
\makeatletter
\renewcommand{\fnum@algorithm}{\fname@algorithm}
\makeatother
\renewcommand{\algorithmicrequire}{\textbf{Input:}}
\renewcommand{\algorithmicensure}{\textbf{Output:}}
\usepackage[margin=1in]{geometry}
\geometry{a4paper}
\usepackage{fancyhdr}
\pagestyle{fancy}
\fancyhf{}
\usepackage{xcolor}

% ------------------------- 代码展示设置 -------------------------
\usepackage{listings}
\lstset{
	basicstyle=\ttfamily,                                % 传统上，代码打字机字体好些
	breaklines,
	tabsize=4,
	extendedchars=false,
	columns=fixed,
	%numbers=left,                                        % 在左侧显示行号
	numbers=none,
	frame=none,                                          % 不显示背景边框
	backgroundcolor=\color[RGB]{245,245,244},            % 设定背景颜色
	keywordstyle=\color[RGB]{40,40,255},                 % 设定关键字颜色
	numberstyle=\footnotesize\color{darkgray},           % 设定行号格式
	commentstyle=\it\color[RGB]{0,96,96},                % 设置代码注释的格式
	stringstyle=\rmfamily\slshape\color[RGB]{128,0,0},   % 设置字符串格式
	showstringspaces=false,                              % 不显示字符串中的空格
	xleftmargin=1em,
	%xrightmargin=1em,
	%aboveskip=-0.5em
	language=c++,                                        % 设置语言
}



% ------------------- 设置超链接，便于跳转 -----------------
\usepackage{enumitem}
\usepackage{hyperref}
\hypersetup{
    pdfborder={0 0 0},    % 消除超链接边框
    colorlinks=true,       % 将边框改为颜色标记（可选）
    linkcolor=black,       % 内部链接颜色（目录、引用等）
    citecolor=black,       % 引用颜色
    urlcolor=blue          % URL链接颜色（保持蓝色以便区分）
}

% ------------------------ 表格设置 ----------------------------
\usepackage{booktabs}
\usepackage{array}
\usepackage{multirow}
\usepackage{longtable}
\usepackage{tabu}
\newcolumntype{L}[1]{>{\raggedright\arraybackslash}p{#1}}
\newcolumntype{M}[1]{>{\centering\arraybackslash}p{#1}}
% ------------------------ 标题区 ----------------------------
\title{集成电路工艺技术基础作业2}
\author{
	姓名：\underline{冯峻}~~~~~~
	学号：\underline{523031910148}~~~~~~}
\date{\today}
\pagenumbering{arabic}

\begin{document}



% ------------------------ 页眉和页脚 ----------------------------
\fancyhead[L]{冯峻}
\fancyhead[C]{集成电路工艺技术基础}
\fancyfoot[C]{\thepage}

\maketitle
%\tableofcontents
%\newpage

\textbf{题目：}

用一张表总结晶圆制造（包括原材料和单晶生长）和清洗的所有流程和目的。

\textbf{解：}

\begin{longtable}{|L{3cm}|L{4.5cm}|L{7cm}|}
\hline
\textbf{阶段} & \textbf{工艺流程} & \textbf{目的与说明} \\
\hline
\endfirsthead

\multicolumn{3}{c}{\textit{续表}} \\
\hline
\textbf{阶段} & \textbf{工艺流程} & \textbf{目的与说明} \\
\hline
\endhead

\hline
\multicolumn{3}{r}{\textit{续下页}} \\
\endfoot

\hline
\endlastfoot

% ============== 原材料制备 ==============
\multirow{3}{*}{\textbf{原材料制备}} 
& 1. 冶金级硅(MGS)制备 & \textbf{原料：}天然砂(SiO$_2$) \\
& & \textbf{反应：}2SiO$_2$ + 2C $\rightarrow$ Si(liquid) + 2CO \\
& & \textbf{纯度：}冶金级，纯度较低，用于后续提纯 \\
\cline{2-3}
& 2. 电子级硅(EGS)提纯 & \textbf{方法：}三氯氢硅法 \\
& & \textbf{反应：}Si + 3HCl $\leftrightarrow$ SiHCl$_3$(liquid) + H$_2$ \\
& & \textbf{纯度：}杂质 < 1 ppb ($\sim$5$\times$10$^{13}$cm$^{-3}$)，O($\sim$10$^{18}$cm$^{-3}$)，C($\sim$10$^{16}$cm$^{-3}$) \\
& & \textbf{用途：}作为单晶生长的原材料 \\
\hline

% ============== 单晶生长 ==============
\multirow{8}{*}{\textbf{单晶生长}} 
& 1. Czochralski(CZ)法 & \textbf{原理：}将EGS硅在SiO$_2$坩埚中加热至1417°C以上熔化，单晶籽晶浸入熔体后缓慢提拉旋转生长 \\
& & \textbf{气氛：}Ar保护气氛 \\
& & \textbf{掺杂：}熔体中加入掺杂剂(如B、P、As等)控制导电类型和电阻率 \\
& & \textbf{控制参数：}提拉速度、熔体温度、旋转速度 \\
& & \textbf{产品：}长晶锭(boule)，质量约100kg，直径可达300-450mm \\
& & \textbf{特点：}商业化主流工艺，O含量高($\sim$10$^{18}$cm$^{-3}$)，C含量适中($\sim$10$^{16}$cm$^{-3}$) \\
& & \textbf{优点：}O提供机械强度和内吸杂(internal gettering)作用 \\
\cline{2-3}
& 2. 浮区(Float Zone, FZ)法 & \textbf{原理：}通过局部加热形成熔区，熔区沿硅棒移动实现区域提纯和单晶生长 \\
& & \textbf{特点：}无坩埚接触，O和C含量极低 \\
& & \textbf{纯度：}比CZ法更高 \\
& & \textbf{应用：}功率器件、聚光太阳能电池 \\
\hline

% ============== 晶圆制造 ==============
\multirow{10}{*}{\textbf{晶圆制造}} 
& 1. 直径研磨 & \textbf{目的：}由于单晶生长过程中直径存在偏差，通过滚磨将晶锭研磨至目标直径 \\
& & \textbf{方法：}机械研磨 \\
\cline{2-3}
& 2. 晶向确认与主平面形成 & \textbf{方法：}X射线衍射确认目标晶向 \\
& & \textbf{目的：}在晶锭上切割参考平面(flat或notch)，标识晶向和掺杂类型 \\
& & \textbf{意义：}便于后续工艺中确定晶体方向 \\
\cline{2-3}
& 3. 切片(Slicing) & \textbf{方法：}使用金刚石切割器将晶锭切成薄片 \\
& & \textbf{目的：}获得厚度约700-800$\mu$m的硅片 \\
& & \textbf{要求：}切割面平整，损伤层最小 \\
\cline{2-3}
& 4. 晶圆抛光 & \textbf{方法：}化学机械抛光(CMP) \\
& & \textbf{目的：}获得无损伤的光滑表面，满足IC制造要求 \\
& & \textbf{控制参数：}曲率、粗糙度 \\
\cline{2-3}
& 5. 标记(Marking) & \textbf{方法：}激光标记 \\
& & \textbf{内容：}晶圆编号、厂商信息等 \\
\cline{2-3}
& 6. 晶圆评估 & \textbf{项目：}晶向检测(通过flat或notch图案识别掺杂类型和晶面方向) \\
& & \textbf{电阻率测试：}四探针法测量电阻率，确认掺杂浓度 \\
& & \textbf{缺陷检测：}检测点缺陷、位错、层错等 \\
\hline

% ============== 晶圆清洗 ==============
\multirow{12}{*}{\textbf{晶圆清洗}} 
& 1. 清洗必要性 & \textbf{原因：}微小污染即可导致器件失效 \\
& & - 移动离子(Na$^+$、K$^+$)：引起阈值电压漂移($\sim$0.1V) \\
& & - 金属杂质(Au、Cu、Fe)：形成深能级陷阱，降低载流子寿命和迁移率 \\
& & - 颗粒污染：导致器件短路、栅氧化层可靠性问题 \\
& & - 有机物：增加栅漏电流，阻碍后续清洗 \\
& & \textbf{要求：}洁净室环境(Class 1-100)和清洗工艺相结合 \\
\cline{2-3}
& 2. 有机物/光刻胶去除 & \textbf{干法清洗：}氧等离子体，将有机物分解为CO$_2$和H$_2$O \\
& & \textbf{湿法清洗：}H$_2$SO$_4$(98\%) : H$_2$O$_2$(30\%) = 2:1-4:1 (Piranha清洗) \\
& & \textbf{温度：}热溶液(通常80-120°C) \\
& & \textbf{目的：}去除有机污染，为后续RCA清洗做准备 \\
\cline{2-3}
& 3. RCA清洗 & \textbf{开发者：}Werner Kern (1965年，RCA公司) \\
& (标准清洗流程) & \textbf{地位：}半导体工艺的标准清洗流程 \\
\cline{2-3}
& 3.1 SC-1清洗 & \textbf{组成：}NH$_4$OH : H$_2$O$_2$ : H$_2$O = 1:1:5 (APM: Ammonia Peroxide Mixture) \\
& (APM) & \textbf{条件：}70-80°C，10分钟 \\
& & \textbf{目的：}去除有机物和颗粒污染 \\
& & \textbf{机理：}轻微腐蚀硅表面，去除颗粒；氧化有机物 \\
\cline{2-3}
& 3.2 稀HF浸泡 & \textbf{组成：}HF : H$_2$O = 1:50 \\
& & \textbf{条件：}室温，1分钟 \\
& & \textbf{目的：}去除SC-1后形成的薄氧化层 \\
\cline{2-3}
& 3.3 SC-2清洗 & \textbf{组成：}HCl : H$_2$O$_2$ : H$_2$O = 1:1:6 (HPM: Hydrochloric Peroxide Mixture) \\
& (HPM) & \textbf{条件：}70-80°C，10分钟 \\
& & \textbf{目的：}去除碱金属离子(Al$^{3+}$、Fe$^{3+}$、Mg$^{2+}$)和重金属(如Au) \\
& & \textbf{机理：}形成可溶性金属氯化物 \\
\cline{2-3}
& 4. 清洗分类 & \textbf{前道清洗(Front-end)：}每次高温工艺前进行，使用强腐蚀性溶液 \\
& & \textbf{后道清洗(Back-end)：}金属布线后使用温和溶液，避免腐蚀金属层 \\
\cline{2-3}
& 5. 洁净室环境 & \textbf{等级：}Class 1-100 (每立方英尺空气中粒径$\geq$0.5$\mu$m的颗粒数) \\
& & \textbf{设施：}HEPA高效过滤器、空气循环系统 \\
& & \textbf{防护：}防尘服("bunny suits")，化学品、水和气体过滤 \\
& & \textbf{目的：}最大限度减少颗粒和污染物 \\
\hline

\caption{晶圆制造与清洗工艺流程总结表}
\label{tab:wafer-process}
\end{longtable}

\textbf{关键知识点补充：}

\begin{enumerate}
\item \textbf{硅晶体缺陷分类}
\begin{itemize}
    \item 点缺陷：替位式、间隙式、空位、肖特基缺陷、弗伦克尔缺陷
    \item 线缺陷：刃位错、螺位错
    \item 面缺陷：层错、晶界、表面
    \item 体缺陷：析出物(Precipitation)
\end{itemize}

\item \textbf{CZ法中O和C的作用}
\begin{itemize}
    \item O来源于SiO$_2$坩埚，典型浓度$\sim$10$^{18}$cm$^{-3}$
    \item C来源于石墨托架，典型浓度$\sim$10$^{16}$cm$^{-3}$
    \item O的优点：提供机械强度；形成SiO$_2$析出物实现内吸杂，捕获有害金属杂质
\end{itemize}

\item \textbf{污染控制要求}
\begin{itemize}
    \item 阈值电压漂移要求：对于10nm栅氧，0.1V漂移对应Q$_M$ = 2.2$\times$10$^{11}$cm$^{-2}$（<0.1\%单层）
    \item 载流子寿命要求：DRAM需要刷新时间>25$\mu$s，要求生成寿命$\tau_G$ > 25$\mu$s
    \item 因此金属杂质(Au、Cu、Fe等)浓度必须极低
\end{itemize}
\end{enumerate}

\end{document}
